\chapter{Carnot Efficiency}

\section*{Introduction}

Imagine two lakes, one hot and one cold. A water wheel extracts energy by letting water flow from the hot lake to the cold lake. The Carnot efficiency says the fraction of heat convertible to work equals the fractional temperature difference. If the hot lake is at 600\,K and the cold at 300\,K, the maximum efficiency is $1 - 300/600 = 50\%$.
\section*{Fundamental Equation Used}

The maximum efficiency of a heat engine between a hot reservoir at $T_H$ and a cold reservoir at $T_C$ is:
\[
\eta_{\text{Carnot}} = 1 - \frac{T_C}{T_H}
\]
where $T_H$ and $T_C$ are absolute temperatures (Kelvin).
\section*{Assumptions}

\textbf{Assumptions:} The engine operates on a reversible cycle (Carnot cycle: two isotherms and two adiabats). The reservoirs are infinite (temperatures do not change). No irreversible processes (no friction, no unrestrained expansion).


\textbf{Limitations:} No real engine is perfectly reversible. The Carnot cycle is infinitely slow (quasi-static), so the power output is zero. Real engines must balance efficiency against power output.


\section*{Mathematical Derivation}

\textbf{Step 1: Define a general heat engine cycle.}

A heat engine operates between a hot reservoir at temperature $T_H$ and a cold reservoir at temperature $T_C$ ($T_H > T_C$). In each cycle, the working substance absorbs heat $Q_H$ from the hot reservoir, rejects heat $Q_C$ to the cold reservoir, and produces work $W = Q_H - Q_C$ (by energy conservation, the first law). The efficiency is:
\begin{equation}
\eta = \frac{W}{Q_H} = \frac{Q_H - Q_C}{Q_H} = 1 - \frac{Q_C}{Q_H}
\end{equation}

\textbf{Step 2: Define the Carnot cycle.}

The Carnot cycle consists of four reversible steps:
\begin{enumerate}
\item \textbf{Isothermal expansion} at $T_H$: the working substance absorbs heat $Q_H$ from the hot reservoir while expanding. Since $\Delta T = 0$, $\Delta U = 0$ (for an ideal gas), so $W_1 = Q_H$.
\item \textbf{Adiabatic expansion}: the substance is insulated and expands, cooling from $T_H$ to $T_C$. No heat exchange ($Q = 0$). Work is done at the expense of internal energy.
\item \textbf{Isothermal compression} at $T_C$: the substance is compressed while rejecting heat $Q_C$ to the cold reservoir. Since $\Delta T = 0$, $\Delta U = 0$, so $W_3 = -Q_C$.
\item \textbf{Adiabatic compression}: the substance is insulated and compressed, heating from $T_C$ back to $T_H$. No heat exchange. Work is done on the substance.
\end{enumerate}

\textbf{Step 3: Apply the second law via entropy.}

For any reversible cycle, the total entropy change is zero:
\begin{equation}
\oint \frac{\delta Q}{T} = 0
\end{equation}
For the four steps of the Carnot cycle, only the isothermal steps exchange heat:
\begin{align}
\Delta S_{\text{isothermal expansion}} &= \frac{Q_H}{T_H} \\
\Delta S_{\text{adiabatic expansion}} &= 0 \\
\Delta S_{\text{isothermal compression}} &= -\frac{Q_C}{T_C} \\
\Delta S_{\text{adiabatic compression}} &= 0
\end{align}
Setting the total to zero:
\begin{equation}
\frac{Q_H}{T_H} - \frac{Q_C}{T_C} = 0
\end{equation}

\textbf{Step 4: Derive the Carnot efficiency.}

From the entropy balance:
\begin{equation}
\frac{Q_C}{Q_H} = \frac{T_C}{T_H}
\end{equation}
Substituting into the efficiency expression:
\begin{equation}
\eta = 1 - \frac{Q_C}{Q_H} = 1 - \frac{T_C}{T_H}
\end{equation}

\textbf{Step 5: Show this is the maximum possible efficiency.}

For any irreversible engine operating between the same reservoirs, the Clausius inequality states:
\begin{equation}
\oint \frac{\delta Q}{T} < 0
\end{equation}
This gives $\frac{Q_H}{T_H} - \frac{Q_C}{T_C} < 0$, so $\frac{Q_C}{Q_H} > \frac{T_C}{T_H}$, and therefore:
\begin{equation}
\boxed{\eta_{\text{Carnot}} = 1 - \frac{T_C}{T_H} \geq \eta_{\text{any engine}}}
\end{equation}
The Carnot efficiency is a universal upper bound, independent of the working substance, determined solely by the reservoir temperatures.

\section*{Characteristics of the Equation}

\begin{itemize}
\item \textbf{Universality:} The efficiency depends only on the temperatures, not on the working substance, engine size, or design details.
\item \textbf{Second Law equivalence:} The Carnot theorem is equivalent to the Kelvin--Planck statement: no process can have the sole result of converting heat entirely into work.
\item \textbf{Reversibility:} Maximum efficiency requires every step to be reversible---the engine passes through equilibrium states.
\item \textbf{Zero power:} A reversible engine takes infinite time per cycle, producing zero power. Real engines sacrifice efficiency for finite power (finite-time thermodynamics).
\item \textbf{Carnot COP:} For heat pumps, the coefficient of performance is $T_H/(T_H - T_C)$ (heating) or $T_C/(T_H - T_C)$ (cooling).
\end{itemize}
\section*{Solution Methods}

\begin{enumerate}
\item \textbf{Carnot cycle analysis:} Follow the working substance through four steps: (1) isothermal expansion at $T_H$ (absorbs $Q_H$), (2) adiabatic expansion (drops to $T_C$), (3) isothermal compression at $T_C$ (rejects $Q_C$), (4) adiabatic compression (rises to $T_H$). Compute $\eta = (Q_H - Q_C)/Q_H$.
\item \textbf{Entropy argument:} For a reversible cycle, total $\Delta S = 0$. Heat absorbed gives $\Delta S = Q_H/T_H$; rejected gives $\Delta S = -Q_C/T_C$. Setting total to zero yields $Q_C/Q_H = T_C/T_H$.
\item \textbf{Clausius inequality:} For any cycle, $\oint \delta Q/T \leq 0$, with equality for reversible cycles. This directly implies the Carnot efficiency is an upper bound.
\end{enumerate}
\section*{Verifications}

\begin{itemize}
\item \textbf{Measured efficiencies:} No real engine has ever exceeded the Carnot efficiency for the same reservoirs, verified across steam engines, gas turbines, and Stirling engines.
\item \textbf{Laboratory demonstrations:} Controlled experiments with quasi-static processes approach but never exceed the Carnot limit.
\item \textbf{Second Law tests:} All attempts to build ``perpetual motion machines of the second kind'' have universally failed.
\end{itemize}
\section*{Applications}

\begin{itemize}
\item \textbf{Power plants:} Steam turbines at $T_H \sim 800$\,K and $T_C \sim 300$\,K give $\eta_{\text{Carnot}} \approx 62.5\%$. Actual efficiencies are $\sim$35--45\%.
\item \textbf{Automotive engines:} Gasoline engines achieve $\sim$25--30\% efficiency, limited by the effective combustion and exhaust temperatures.
\item \textbf{Refrigeration:} The Carnot COP limits cooling systems, guiding development of efficient refrigerants and compressors.
\item \textbf{Solar energy:} The Sun ($\sim$5800\,K) and Earth ($\sim$300\,K) give $\eta_{\text{Carnot}} \approx 95\%$, though practical solar cells achieve $\sim$20--25\%.
\item \textbf{Thermoelectrics:} The Carnot limit guides development of thermoelectric generators and coolers.
\end{itemize}
\section*{What Richard Feynman Would Have Asked}

\subsection*{1. Plain-English Translation}
\textit{``What does the Carnot efficiency really mean for someone who builds engines?''}

The Carnot efficiency is nature's speed limit for converting heat into work. It says: no matter how good your engineer is, you cannot extract more than a certain fraction of heat energy as useful work. The rest must be dumped into a cold reservoir. This is not because our engines are poorly designed---it is because heat has a directional quality. Heat flows from hot to cold, and you can only extract work from that flow. The Carnot efficiency tells you exactly how much work is hiding in a given temperature difference.

The practical implication: to get more work from a heat engine, you must either raise the hot temperature or lower the cold temperature. This is why jet engines run hot (turbine inlet temperatures above 1500\,K) and why power plants use cooling towers or river water as cold sinks.

\subsection*{2. Localized Mechanics}
\textit{``What is happening to the molecules inside the engine during each step of the Carnot cycle?''}

During isothermal expansion at $T_H$, molecules absorb heat from the hot reservoir. Their average kinetic energy stays constant, but they do work by expanding---pushing on a piston. The energy comes from the heat flow, not from the molecules speeding up. During adiabatic expansion, no heat flows, so the molecules do work at the expense of their own kinetic energy---they slow down, and the temperature drops.

He would press you on the entropy change: during isothermal expansion, entropy increases (heat flows in). During isothermal compression, entropy decreases (heat flows out). During adiabatic steps, entropy is constant. Over the full cycle, total entropy change is zero---the engine returns to its initial state. This zero net entropy change is the mathematical expression of reversibility and the reason the Carnot engine achieves maximum efficiency.

\subsection*{3. Testing Extreme Limits}
\textit{``What happens when we push the temperatures to extreme values?''}

As $T_C \to 0$ (approaching absolute zero), $\eta_{\text{Carnot}} \to 1$---converting all heat to work. But the Third Law says you cannot reach absolute zero in finite steps, so this limit is unattainable. As $T_H \to \infty$, efficiency also approaches 1, but materials melt long before extremely high temperatures.

The most extreme known heat engines are stellar: the Sun operates between $\sim$5800\,K (surface) and $\sim$3\,K (cosmic microwave background), giving $\eta_{\text{Carnot}} \approx 99.95\%$. Stars are incredibly efficient heat engines, converting nuclear energy to radiation with very little waste.

\subsection*{4. Dimensionality and Geometry}
\textit{``Does the shape or size of the engine affect the Carnot efficiency?''}

No. The Carnot efficiency depends only on the two temperatures---not on the size, shape, pistons, working substance, or any design detail. This universality is remarkable: a tiny laboratory engine and a massive power plant, operating between the same temperatures, have the same maximum efficiency. Geometry affects power output (a bigger engine produces more per cycle), but not efficiency.

\subsection*{5. Cross-Disciplinary Connection}
\textit{``Where else does this same efficiency bound appear outside of heat engines?''}

The Carnot limit has deep analogues in information theory. Landauer's principle states that erasing one bit of information requires at least $kT \ln 2$ of energy dissipated as heat---an ``information engine'' faces the same Carnot-like constraints. The maximum efficiency of any process converting free-energy difference into work is bounded by the ratio of reservoir temperatures. In biology, molecular motors (kinesin, ATP synthase) operate under similar thermodynamic constraints, with efficiencies limited by the chemical potential differences of their fuel molecules.


%=============================================================
% Chapter 16: Clausius-Clapeyron Equation
%=============================================================
\section*{FAQ}

\textbf{Q1: Can we reach the Carnot efficiency in practice?}

A: Only in principle, by operating infinitely slowly through equilibrium states. Real engines have irreversibilities that lower efficiency. The goal is to minimize these and approach the Carnot limit as closely as practical.

\textbf{Q2: Why does efficiency depend on temperature ratio rather than difference?}

A: Because thermodynamics is about ratios. The entropy exchanged as $Q/T$ must balance, and $T_C/T_H$ emerges from this balance. A 100\,K difference between 300\,K and 400\,K gives 25\% efficiency; between 1000\,K and 1100\,K, only 10\%.

\textbf{Q3: Is the Carnot efficiency a law or just an idealization?}

A: It is a provable upper bound---a law of physics. It follows from the Second Law, tested for over 150 years with no exceptions. The bound is tight: reversible engines can approach it arbitrarily closely.
\section*{Important Bibliography}

\begin{enumerate}
\item Carnot, S., \textit{R\'eflexions sur la puissance motrice du feu} (Paris, 1824); English trans. by R.H. Thurston (Dover, 1960).
\item Planck, M., \textit{Treatise on Thermodynamics}, 3rd ed. (Dover, 1945), Chapter 170.
\item Callen, H.B., \textit{Thermodynamics and an Introduction to Thermostatistics}, 2nd ed. (Wiley, 1985), Chapter 4.
\item Fermi, E., \textit{Thermodynamics} (Dover, 1956), Chapter III.
\end{enumerate}


